\chapter{Diffusion Models} \label{ch:diffusion}
\index{Diffusion Models}

In recent years, generative models have seen remarkable success, with diffusion models emerging as the state-of-the-art technique for high-fidelity image generation, superseding earlier approaches such as Generative Adversarial Networks (GANs\index{GANs}). Unlike GANs, which generate images in a single forward pass mapping a latent vector $\mathbf{z} \sim \phi_z$ to an image, or autoregressive models that generate sequences token-by-token, diffusion models define a Markov chain of diffusion steps to slowly add random noise to data and then learn to reverse the diffusion process to construct desired data samples from the noise.

\section{The Forward Diffusion Process} \label{sec:diffusion_forward}
\index{Forward Diffusion Process}

The core idea of diffusion models is based on non-equilibrium thermodynamics. Given a data point from a real data distribution $\mathbf{x}_0 \sim q(\mathbf{x})$, the forward diffusion process adds Gaussian noise in $T$ steps, producing a sequence of noisy samples $\mathbf{x}_1, \mathbf{x}_2, \dots, \mathbf{x}_T$. The forward process is defined as a Markov chain with Gaussian transitions:
\begin{equation} \label{eq:forward_transition}
    q(\mathbf{x}_t | \mathbf{x}_{t-1}) = \mathcal{N}\left(\mathbf{x}_t; \sqrt{1 - \beta_t} \mathbf{x}_{t-1}, \beta_t \mathbf{I}\right)
\end{equation}
where the step sizes are controlled by a pre-defined variance schedule $\{\beta_t \in (0, 1)\}_{t=1}^T$. Given a sufficiently large $T$ and a well-behaved schedule of $\beta_t$, the final latent variable $\mathbf{x}_T$ converges to an isotropic Gaussian distribution $\mathcal{N}(\mathbf{0}, \mathbf{I})$.

A remarkably ``nice property'' of the forward process is that we can sample $\mathbf{x}_t$ at any arbitrary time step $t$ in a closed form using the reparameterization trick\index{Reparameterization Trick}. Let $\alpha_t = 1 - \beta_t$ and $\bar{\alpha}_t = \prod_{i=1}^t \alpha_i$. We can rewrite the formulation as:
\begin{equation} \label{eq:forward_closed_form}
    \mathbf{x}_t = \sqrt{\alpha_t} \mathbf{x}_{t-1} + \sqrt{1 - \alpha_t} \boldsymbol{\epsilon}_{t-1} = \dots = \sqrt{\bar{\alpha}_t} \mathbf{x}_0 + \sqrt{1 - \bar{\alpha}_t} \boldsymbol{\epsilon}
\end{equation}
where $\boldsymbol{\epsilon} \sim \mathcal{N}(\mathbf{0}, \mathbf{I})$. Consequently, the marginal distribution of $\mathbf{x}_t$ conditioned on $\mathbf{x}_0$ is given by:
\begin{equation} \label{eq:forward_marginal}
    q(\mathbf{x}_t | \mathbf{x}_0) = \mathcal{N}\left(\mathbf{x}_t; \sqrt{\bar{\alpha}_t} \mathbf{x}_0, (1 - \bar{\alpha}_t) \mathbf{I}\right)
\end{equation}

\begin{insightbox}[Intuition: The Forward Process]
Think of the forward diffusion process as a slow deterioration of a pristine image into complete static, like a television losing its signal. At each discrete step $t$, we inject a tiny, carefully calibrated amount of Gaussian noise. By the time we reach $T$ (often $T=1000$), all the original semantic and perceptual information is entirely destroyed, leaving us with pure, unstructured noise. The closed-form equation is powerful because it means we do not need to iterate through all previous steps $1$ to $t-1$ to find out what the noisy image looks like at step $t$; we can jump there instantly by adding scaled noise to the original image $\mathbf{x}_0$.
\end{insightbox}

\section{The Reverse Diffusion Process} \label{sec:diffusion_reverse}
\index{Reverse Diffusion Process}

If we could reverse the forward process and sample from $q(\mathbf{x}_{t-1} | \mathbf{x}_t)$, we could start from pure Gaussian noise $\mathbf{x}_T \sim \mathcal{N}(\mathbf{0}, \mathbf{I})$ and recreate the true sample from the real data distribution. However, $q(\mathbf{x}_{t-1} | \mathbf{x}_t)$ depends on the entire data distribution and is intractable to compute. Thus, we learn a parameterized model $p_\theta(\mathbf{x}_{t-1} | \mathbf{x}_t)$ to approximate it. If $\beta_t$ is sufficiently small, the reverse transitions are also Gaussian:
\begin{equation} \label{eq:reverse_transition}
    p_\theta(\mathbf{x}_{t-1} | \mathbf{x}_t) = \mathcal{N}\left(\mathbf{x}_{t-1}; \boldsymbol{\mu}_\theta(\mathbf{x}_t, t), \boldsymbol{\Sigma}_\theta(\mathbf{x}_t, t)\right)
\end{equation}
Common formulations simplify this by keeping the variance fixed, setting $\boldsymbol{\Sigma}_\theta(\mathbf{x}_t, t) = \sigma_t^2 \mathbf{I}$ (where $\sigma_t^2 = \beta_t$ or a similar related constant), meaning the neural network only needs to learn the mean $\boldsymbol{\mu}_\theta(\mathbf{x}_t, t)$.

\begin{insightbox}[Intuition: The Reverse Process]
The reverse process is the actual ``generation'' phase. We start with pure static (Gaussian noise) and ask the neural network to repeatedly denoise the image step by step. The model does not predict the pristine image $\mathbf{x}_0$ in one shot; rather, it predicts the exact noise vector that was added at timestep $t$. We then subtract a tiny fraction of this predicted noise, gradually carving out an image from the static. Over hundreds of steps, a coherent image emerges.
\end{insightbox}

\section{Training and Inference} \label{sec:diffusion_training}
\index{Diffusion Models!Training}
\index{Diffusion Models!Inference}

To train the reverse diffusion process, we aim to minimize the difference between the true posterior $q(\mathbf{x}_{t-1} | \mathbf{x}_t, \mathbf{x}_0)$ and the learned distribution $p_\theta(\mathbf{x}_{t-1} | \mathbf{x}_t)$. It is known that the true posterior mean is given by:
\begin{equation} \label{eq:true_posterior_mean}
    \tilde{\boldsymbol{\mu}}_t(\mathbf{x}_t, \mathbf{x}_0) = \frac{1}{\sqrt{\alpha_t}} \left( \mathbf{x}_t - \frac{1 - \alpha_t}{\sqrt{1 - \bar{\alpha}_t}} \boldsymbol{\epsilon}_t \right)
\end{equation}
Therefore, we train the neural network to predict $\boldsymbol{\epsilon}_t$, the noise added to the image at time $t$. The reparameterized mean predicted by our model becomes:
\begin{equation} \label{eq:predicted_mean}
    \boldsymbol{\mu}_\theta(\mathbf{x}_t, t) = \frac{1}{\sqrt{\alpha_t}} \left( \mathbf{x}_t - \frac{1 - \alpha_t}{\sqrt{1 - \bar{\alpha}_t}} \boldsymbol{\epsilon}_\theta(\mathbf{x}_t, t) \right)
\end{equation}
Ho et al. (2020) demonstrated empirically that the model works better when trained with a simplified objective function that drops the complex weighting terms. The simplified loss function $L_t^{\text{simple}}$ minimizes the mean squared error between the actual noise $\boldsymbol{\epsilon}_t$ and the predicted noise $\boldsymbol{\epsilon}_\theta$:
\begin{equation} \label{eq:simple_loss}
    L_t^{\text{simple}} = \mathbb{E}_{t \sim [1,T], \mathbf{x}_0, \boldsymbol{\epsilon}_t} \left[ \left\| \boldsymbol{\epsilon}_t - \boldsymbol{\epsilon}_\theta(\mathbf{x}_t, t) \right\|^2 \right] = \mathbb{E}_{t \sim [1,T], \mathbf{x}_0, \boldsymbol{\epsilon}_t} \left[ \left\| \boldsymbol{\epsilon}_t - \boldsymbol{\epsilon}_\theta\left(\sqrt{\bar{\alpha}_t} \mathbf{x}_0 + \sqrt{1 - \bar{\alpha}_t} \boldsymbol{\epsilon}_t, t\right) \right\|^2 \right]
\end{equation}

During \textbf{inference} (sampling), we begin with $\mathbf{x}_T \sim \mathcal{N}(\mathbf{0}, \mathbf{I})$. For each step from $t=T$ down to $1$, we sample $\mathbf{z} \sim \mathcal{N}(\mathbf{0}, \mathbf{I})$ (if $t>1$, else $\mathbf{z}=\mathbf{0}$) and apply the denoising update:
\begin{equation} \label{eq:sampling_update}
    \mathbf{x}_{t-1} = \frac{1}{\sqrt{\alpha_t}} \left( \mathbf{x}_t - \frac{1 - \alpha_t}{\sqrt{1 - \bar{\alpha}_t}} \boldsymbol{\epsilon}_\theta(\mathbf{x}_t, t) \right) + \sigma_t \mathbf{z}
\end{equation}

\subsection{Variance Scheduling} \label{subsec:scheduling}
\index{Variance Scheduling}
\index{Linear Scheduling}
\index{Cosine Scheduling}

The choice of the variance schedule $\beta_t$ plays a critical role. Early works such as Ho et al. (2020) used a \textit{linear schedule}, where $\beta_t$ increases linearly from $10^{-4}$ to $0.02$. However, linear scheduling destroys the information too quickly, bringing the latent variable to pure noise prematurely.

Nichol \& Dhariwal (2021) proposed a \textit{cosine schedule}, defined by:
\begin{equation} \label{eq:cosine_schedule}
    \bar{\alpha}_t = \frac{f(t)}{f(0)}, \quad \text{where } f(t) = \cos^2\left( \frac{t/T + s}{1 + s} \cdot \frac{\pi}{2} \right)
\end{equation}
The small offset $s$ prevents $\beta_t$ from becoming exceedingly small near $t=0$. This schedule ensures a more gradual transition towards pure noise, improving the overall quality and log-likelihood of generated samples.

\section{Guidance in Diffusion Models} \label{sec:diffusion_guidance}
\index{Guidance}

Generating random images is useful, but we often want to condition the generation on class labels or textual descriptions. Two primary methods for conditional generation have emerged.

\subsection{Classifier Guided Diffusion}
\index{Classifier Guided Diffusion}

Dhariwal \& Nichol (2021) introduced Classifier Guided Diffusion. By training an explicit classifier $f_\phi(y|\mathbf{x}_t)$ on noisy images, we can utilize its gradients to guide the diffusion sampling process toward the target condition $y$. The score function of the joint distribution modifies the noise predictor:
\begin{equation} \label{eq:classifier_guidance}
    \bar{\boldsymbol{\epsilon}}_\theta(\mathbf{x}_t, t) = \boldsymbol{\epsilon}_\theta(\mathbf{x}_t, t) - w \sqrt{1 - \bar{\alpha}_t} \nabla_{\mathbf{x}_t} \log f_\phi(y|\mathbf{x}_t)
\end{equation}
where $w$ is the gradient scale. This method allowed the Ablated Diffusion Model with Guidance (ADM-G) to achieve better results than state-of-the-art generative models like BigGAN.

\subsection{Classifier-Free Guidance}
\index{Classifier-Free Guidance}

Alternatively, Ho \& Salimans (2021) introduced \textit{Classifier-Free Guidance} (CFG), which eliminates the need for an external classifier. In this approach, a single neural network is jointly trained for both unconditional and conditional generation. Let $p_\theta(\mathbf{x}|\emptyset)$ be parameterized through a score estimator $\boldsymbol{\epsilon}_\theta(\mathbf{x}_t, t)$ and $p_\theta(\mathbf{x}|y)$ through $\boldsymbol{\epsilon}_\theta(\mathbf{x}_t, t, y)$. During training, the conditioning signal $y$ is randomly replaced with a learnable null token $\emptyset$ with some probability. 

During sampling, the model makes two predictions: one conditional and one unconditional. The final extrapolated noise prediction is:
\begin{equation} \label{eq:cfg}
    \bar{\boldsymbol{\epsilon}}_\theta(\mathbf{x}_t, t, y) = (w + 1)\boldsymbol{\epsilon}_\theta(\mathbf{x}_t, t, y) - w \boldsymbol{\epsilon}_\theta(\mathbf{x}_t, t, \emptyset)
\end{equation}
This pushes the image generation away from generic, unconditional priors and strongly toward the conditioning prompt.

\section{Improving Sampling Efficiency} \label{sec:diffusion_sampling}

A fundamental drawback of standard DDPMs is the requirement for hundreds or thousands of forward passes during inference, making sampling slow. Several solutions have been proposed:

\begin{itemize}
    \item \textbf{Progressive Distillation\index{Progressive Distillation}}: Salimans \& Ho (2022) proposed distilling a trained deterministic sampler into a new student model that halves the required sampling steps at each distillation iteration. A student DDIM\index{DDIM} step is trained to match two teacher steps.
    \item \textbf{Consistency Models\index{Consistency Models}}: Song et al. (2023) introduced consistency models that learn to map any intermediate noisy data point $\mathbf{x}_t$ directly back to its origin $\mathbf{x}_0$. Given a consistency function $f$, it enforces $f(\mathbf{x}_t, t) = f(\mathbf{x}_{t'}, t') = \mathbf{x}_\epsilon$ for any $t, t'$. This allows trading computation for quality, enabling high-quality single-step generation.
\end{itemize}

\section{Scaling Architectures: LDM and DiT} \label{sec:diffusion_architectures}

Historically, diffusion models relied heavily on the \textbf{U-Net\index{U-Net}} backbone architecture, originally developed for image segmentation. The U-Net employs a sequence of downsampling blocks, self-attention blocks, and upsampling blocks with skip connections. To incorporate temporal information, a sinusoidal time embedding is added to all blocks. For text conditioning, an embedder integrates features via cross-attention mechanisms.

\subsection{Latent Diffusion Models (LDM)}
\index{Latent Diffusion Models (LDM)}

Operating directly in the high-dimensional pixel space is computationally prohibitive. Rombach \& Blattmann et al. (2022) introduced Latent Diffusion Models (LDMs). Their insight is that most bits of an image contribute merely to high-frequency perceptual details, while the semantic composition remains intact under compression. 

LDMs use an Autoencoder\index{Autoencoder} (an encoder $\mathcal{E}$ and decoder $\mathcal{D}$) to compress the input image $\mathbf{x} \in \mathbb{R}^{H \times W \times 3}$ into a smaller 2D latent vector $\mathbf{z} = \mathcal{E}(\mathbf{x}) \in \mathbb{R}^{h \times w \times c}$, trimming off pixel-level redundancy. The diffusion process is then executed efficiently within this low-dimensional latent space. Regularization methods such as KL-penalty (KL-reg\index{KL-reg}) or vector quantization (VQ-reg\index{VQ-reg}) are applied during the autoencoder training to avoid arbitrarily high-variance in latent spaces.

\subsection{Diffusion Transformers (DiT)}
\index{Diffusion Transformer (DiT)}

Moving beyond the U-Net paradigm, Peebles \& Xie (2023) proposed the Diffusion Transformer (DiT). DiT operates on the latent patches using the same design space of the Latent Diffusion Model but utilizes a Transformer backbone. The latent representation $\mathbf{z}$ is ``patchified'' into a sequence of patches of size $p$, converting it into a sequence of $(I/p)^2$ patches. This sequence of tokens is processed through standard Transformer blocks, using adaptive layer normalization (adaLN) for conditioning on the timestep $t$ and label $y$. The transformer decoder outputs the required noise predictions and an output diagonal covariance prediction, scaling effectively with network depth and width.
